\documentclass[11pt,a4paper]{article}
\usepackage{operationsform}
\begin{document}
\OPStitle{取材前確認チェックリスト}{TPL-CHK-021}{0.2.0}
\OPSsection{1}{確認項目}
\OPScheck{1}{取材目的、想定媒体、読者、取材者の役割を文書で確認した}
\OPScheck{2}{調査資料と主張IDを読み、何を検証する取材か説明できる}
\OPScheck{3}{対象者の選定理由と必要な別視点を記録した}
\OPScheck{4}{質問案、実施予定、時間配分、連絡方法を共有した}
\OPScheck{5}{録音・撮影・氏名掲載・公開前確認の未確定条件を抽出した}
\OPScheck{6}{センシティブな話題、機材、アクセス配慮、安全対策を検討した}
\OPScheck{7}{担当者と責任者が当日の準備状況と未決項目を確認した}
\OPScheck{8}{適用外・実施不能項目について理由を記録した}
\OPSsection{2}{点検結果の記録}
\OPSblankshort{対象案件・対象版}
\OPSblankshort{担当者と点検日}
\OPSblanklong{未解決事項と対応}
\par\Needspace{4\baselineskip}\noindent\textbf{確認に用いた資料ID}\par\vspace{1mm}
{\fontsize{8.5}{13}\selectfont
\begin{tabularx}{\linewidth}{@{}XXXX@{}}
\textbf{対象・項目} & \textbf{内容・参照} & \textbf{担当・状態} & \textbf{備考}\\[2mm]
\strut & \strut & \strut & \strut\\[5mm]
\strut & \strut & \strut & \strut\\[5mm]
\end{tabularx}
}\vspace{2mm}\par
\end{document}
